\label{ch:zustand_prozess}
% Quelldokument: 247

\epigraph{Die bloße Speicherung von Information ist energieneutral.
Erst die Verarbeitung hat eine Energiedimension.}{Dok.\,247}

\section{Die Frage}

In der Informationsphysik --- vertreten durch Wheeler (``It from Bit''),
M.\,M.~Vopson \textbf{[V1--V3]} und andere --- wird Information oft als ontologisches Primitiv
behandelt, aus dem Energie und Materie folgen. Landauers Prinzip wird dabei
als Beleg angeführt: Da das Löschen eines Bits Energie kostet,
sei Information physikalisch real.

Dieser Schluss verwechselt zwei strukturell verschiedene Begriffe.

\begin{keyresult}
\textbf{Information als Zustand} --- eine Konfiguration, ein gespeichertes Bit,
eine Wicklungszahl --- existiert ohne Energieaufwand. Sie ist eine topologische
oder geometrische Eigenschaft.\\[4pt]
\textbf{Information als Prozess} --- Lesen, Schreiben, Löschen, Transformieren
--- kostet Energie. Landauers Prinzip gilt \emph{nur} für die Löschung, also
für einen Prozess, nicht für den Zustand selbst.
\end{keyresult}

\section{Empirischer Beleg: der Permanentspeicher}

Wer den Unterschied einebnet, muss eine direkte empirische Konsequenz
akzeptieren: Jedes gespeicherte Bit auf einem ausgeschalteten Gerät würde
kontinuierlich Energie verbrauchen --- was Permanentspeicher grundsätzlich
unmöglich machen würde. Die Existenz von Flash-Speicher, magnetischen Festplatten
und ROM-Bausteinen widerlegt die Gleichsetzung von Information und Energie.

Ein Permanentspeicher hält Millionen von Bits über Jahre bei nahezu null
Energieverbrauch. Er ist ein Gegenexperiment zu jeder Theorie, die Information
und Energie gleichsetzt.

\section{Die FFGFT-Lesart}

In FFGFT sind Zustände geometrische Konfigurationen auf $T^4$:
Wicklungszahlen $(n_\theta, n_\varphi) \in \mathbb{Z}^2$. Energie ist eine
abgeleitete Größe, die aus \emph{Prozessen} folgt --- aus dem Übergang zwischen
Konfigurationen --- nicht aus der bloßen Existenz einer Konfiguration.

Die Ruhemasse eines Teilchens in FFGFT ist:
\[
  m = \frac{\hbar c}{L} \cdot n \qquad (n \in \mathbb{Z},\; L \text{ Skala})
\]
Die Wicklungszahl $n$ ist ein Zustand --- sie kostet keine Energie, solange
die Topologie intakt bleibt. Erst der Übergang $n \to n'$ --- ein Prozess ---
erfordert oder erzeugt Energie.

\section{M.\,M.~Vopson und die FFGFT-Einschränkung}

M.\,M.~Vopson postuliert in \textbf{[V1]} eine Massen-Information-Äquivalenz: Jedes Bit
trägt eine universelle Ruhemasse. Das ist attraktiv --- und nach FFGFT falsch.

FFGFT-Einschränkung: Vopsons Bitmassse wäre systeminvariant. FFGFT zeigt,
dass die Energie eines Bits $\Ebit = \hbar c / L$ von der Skala $L$ abhängt
(Dok.\,257). Bits auf verschiedenen Skalen haben verschiedene Energien.
Eine universelle Bitmasse gibt es nicht --- weder nach Landauer (Dok.\,A271--A272)
noch nach FFGFT.

\begin{ffgftbox}
\textbf{Wheeler:} ``It from Bit'' --- Geometrie aus Information.\\
\textbf{M.\,M.~Vopson \textbf{[V1,V2]}:} Bitmasse universell --- Information hat Ruhemasse.\\
\textbf{FFGFT:} ``Bit from It'' --- Information als Projektion aus Geometrie;
Bitmasse skalensensitiv; Speichern ist kostenlos, Löschen kostet.
\end{ffgftbox}

\section{Was Landauer wirklich zeigt}

Landauers Prinzip ist in dieser Perspektive kein Beleg für die
physikalische Realität von Information als ontologischem Primitiv.
Es ist ein Beleg für die physikalische Realität des \emph{Trägers}:
Wenn ein Träger irreversibel umkonfiguriert wird, fallen Kosten an.
Die Information auf dem Träger spielt dabei keine Rolle --- die
Thermodynamik zählt Gebiete, nicht Bedeutungen (Dok.\,A271).
